% !PRINTED 572
\weberpage{572}

%<<CONT
\weq{1}{
\begin{aligned}
&\text{a)} \quad \tau(u) = \frac{g_2 g_3}{G} \wp(u), \quad \text{im allgemeinen,} \\
&\text{b)} \quad \phantom{\tau(u)} = \frac{\wp(u)^2}{g_2}, \quad \text{wenn } g_3 = 0, \\
&\text{c)} \quad \phantom{\tau(u)} = \frac{\wp(u)^3}{g_3}, \quad \text{wenn } g_2 = 0.
\end{aligned}
}

Und wenn dann $\mu$ wie oben eine ganze Zahl des Körpers $\Omega$ ist, so können wir in allen drei Fällen setzen:

\weq{2}{\tau(\mu u) = \frac{\Psi(\tau)}{\Phi(\tau)},}

worin $\Phi, \Psi$ ganze Funktionen von $\tau$ sind, die bis auf konstante Faktoren im Falle (1) a) mit $R, P$ des vorigen Paragraphen übereinstimmen, in den Fällen (1) b) und (1) c) daraus durch Erhebung ins Quadrat oder in den Kubus hervorgehen [§\,152, 4), 5)].

Wir richten den Bruch (2) so ein, daß die höchste Potenz von $\tau$ im Nenner $\Phi$ den Koeffizienten 1 hat. Die übrigen Koeffizienten von $\Phi$ sind dann rational zusammengesetzt aus Größen der Form

\[
\tau\left(\frac{\nu}{\mu}\right) = \tau\left(\frac{\nu \mu'}{m}\right),
\]

und da $\nu\mu'$ eine Periode von $\tau$ ist, so gehören diese Größen zu den Wurzeln der Teilungsgleichung der Perioden, und sind daher, da der Modul der entsprechenden elliptischen Funktionen eine algebraische Zahl ist, selbst algebraische Zahlen (§\,58, §\,61).

Daraus folgt:

\begin{quote}
2. \gesperrt{Die Koeffizienten in $\Phi(\tau)$ sind algebraische Zahlen.}
\end{quote}

Wir erweitern den Bruch $\Psi : \Phi$ durch Multiplikation von Zähler und Nenner mit dem Produkt der konjugierten Werte $\Phi', \Phi'', \Phi''', \ldots$ und erhalten:

\weq{3}{\tau(\mu u) = \frac{Z}{N},}

worin nun $Z$ und $N$ ganze Funktionen von $\tau$ sind, die einen gemeinsamen Teiler enthalten können, wobei jedoch $N$ rationale Zahlenkoeffizienten hat.
